\section{Corrected S-tier-superlevel cut and the affine151
branch}

\textbf{Status:} publication dependency replacing the cut/count portions
of \texttt{global\_atlas\_interface\_closure.md} and
\texttt{stoichiometric\_gate\_feasibility.md}. The audited legacy files
are retained unchanged for provenance.

\textbf{Date:} 2026-08-12 PDT.

\textbf{Retirement declaration.} Section 8.2 of the frozen
\texttt{global\_atlas\_interface\_closure.md} and the old stoichiometric
incidence totals \(12886/9913/2973\) are retired as publication
dependencies. The corrected totals are \(12678/9709/2969\). The
tier-certified and residual pair fingerprints, the exact 2,511 pair
residual, and the affine151 recurrence branch remain unchanged and are
repinned below.

\subsection{1. Exact scope}

Consider the finite binary complex universe

\[
 \mathcal C_2=\{0,A,B,C,2A,2B,2C,A+B,A+C,B+C\}
\]

and an ordered pair of disjoint linkage supports, each equipped with an
arbitrary strongly connected directed reaction graph and arbitrary
positive rates. This note supplies:

\begin{enumerate}
\def\labelenumi{\arabic{enumi}.}
\tightlist
\item
  the exact symbolic condition under which \textbf{every} such pair of
  strongly connected graphs has an Anderson--Kim descending source for a
  fixed tier descriptor;
\item
  the exact finite pair split after the pre-tier support branches;
\item
  the corrected affine-feasibility incidence totals; and
\item
  the class-local recurrence conclusion for the unchanged 151-pair set.
\end{enumerate}

There is no enumeration of orientations, stochastic paths, population
states, or rate values. Finite computation only deduplicates rational
tier types and checks support/tier/affine-subspace identities.

\subsection{2. Exact universal strong-orientation
theorem}

Fix a tier descriptor. Write \(E=T^{S,1}\) for the global top S-tier
among the network complexes. If \(E\ne\varnothing\), all its complexes
occupy a single D-tier; call its level \(r\). For each linkage support
\(L\), define

\[
 U_L(r)=\{y\in L:y\text{ lies at or above D-level }r\}.
 \tag{2.1}
\]

\begin{quote}
\textbf{Theorem 2.1 (S-tier-superlevel cut).} Every strongly connected
directed graph on each linkage support has a D-descending reaction whose
source lies in \(E\) if and only if, for some linkage \(L\), \[
  \varnothing\ne U_L(r)\subsetneq L,
  \qquad U_L(r)\subseteq E.
  \tag{2.2}
\]
\end{quote}

\subsubsection{Proof}

If (2.2) holds, strong connectivity gives a directed path from
\(U_L(r)\) to its complement. The first exiting edge has source in
\(U_L(r)\subseteq E\) and target below D-level \(r\). It is the required
descending reaction.

Conversely, suppose (2.2) fails. Work independently in each linkage. If
\(U_L(r)\) is empty or all of \(L\), choose any directed Hamiltonian
cycle. Otherwise choose \(b\in U_L(r)\setminus E\), which exists by
failure of (2.2), and order a directed Hamiltonian cycle as

\[
 U_L(r)\setminus\{b\},\quad b,\quad L\setminus U_L(r).
 \tag{2.3}
\]

Its only edge from \(U_L(r)\) to its complement is sourced at
\(b\notin E\). Every vertex of \(E\cap L\) has its successor within
\(U_L(r)\), where no target is below level \(r\). Hence this strongly
connected cycle has no E-sourced descending edge. Combining the cycles
proves the converse. \(\square\)

The superlevel in (2.1), not merely the global top D-tier, is essential.
For

\[
 L_1=\{C,2C\},\quad L_2=\{A,2A,A+B,A+C\},\quad
 w=(0,3,1),\quad c=(0,2,2),
\]

the global top D-tier is the disabled singleton \(\{A+B\}\), whereas the
top S-tier is \(E=\{2C\}\). Here \(U_{L_1}=\{2C\}\), so every strongly
connected graph on \(L_1\) supplies a \(2C\)-sourced descent. A test
using only the global top D-tier incorrectly rejects this descriptor.

\subsection{3. Completeness of the finite tier
list}

For a tier sequence put

\[
 u_n=(\log(x_{n,A}\vee1),\log(x_{n,B}\vee1),
      \log(x_{n,C}\vee1)).
\]

Project away the span of every D-equivalent complex difference and every
coordinate vector belonging to an eventually bounded coordinate. The
projection removed from \(u_n\) is bounded. Strict D-comparisons
therefore remain strict, D-equalities become exact equalities, active
coordinates stay positive, and bounded coordinates become zero. The
resulting nonempty relatively open rational cone contains a rational
weight. Thus an ordinary nonnegative rational weight represents every
D-preorder, including multiscale preorders; no polynomial relation among
the original population scales is assumed.

The 21 projectively distinct comparison hyperplanes on the normalized
nonnegative simplex have 37 vertices. Vertices, edge midpoints, and
triangle centroids give 5,128 rational candidates and 193 distinct
\((\text{D-preorder},\text{active-coordinate set})\) types. An
eventually bounded integer coordinate is eventually constant, and
binary-source availability depends only on whether that constant is
\(0\), \(1\), or at least \(2\). Adding these exact caps yields 259
descriptors.

Theorem 2.1 reduces the universal quantifier over strongly connected
graphs to the finite set containment (2.2). The program evaluates that
containment only; it does not construct or enumerate orientations.

\subsection{4. Pre-tier support branches and the exact 2,511
residual}

The ordered support branches are applied in this order: finite strict
invariant, active-chart invariant, full deficiency zero, exact
signed-service seam, exact residual pair, exact seven-support seam, then
tier certification.

{\small\def\LTcaptype{none} % do not increment counter
\begin{longtable}[]{@{}
  >{\raggedright\arraybackslash}p{(\linewidth - 16\tabcolsep) * \real{0.0857}}
  >{\raggedleft\arraybackslash}p{(\linewidth - 16\tabcolsep) * \real{0.1143}}
  >{\raggedleft\arraybackslash}p{(\linewidth - 16\tabcolsep) * \real{0.1143}}
  >{\raggedleft\arraybackslash}p{(\linewidth - 16\tabcolsep) * \real{0.1143}}
  >{\raggedleft\arraybackslash}p{(\linewidth - 16\tabcolsep) * \real{0.1143}}
  >{\raggedleft\arraybackslash}p{(\linewidth - 16\tabcolsep) * \real{0.1143}}
  >{\raggedleft\arraybackslash}p{(\linewidth - 16\tabcolsep) * \real{0.1143}}
  >{\raggedleft\arraybackslash}p{(\linewidth - 16\tabcolsep) * \real{0.1143}}
  >{\raggedleft\arraybackslash}p{(\linewidth - 16\tabcolsep) * \real{0.1143}}@{}}
\toprule\noalign{}
\begin{minipage}[b]{\linewidth}\raggedright
family
\end{minipage} & \begin{minipage}[b]{\linewidth}\raggedleft
unique pairs
\end{minipage} & \begin{minipage}[b]{\linewidth}\raggedleft
finite
\end{minipage} & \begin{minipage}[b]{\linewidth}\raggedleft
active chart
\end{minipage} & \begin{minipage}[b]{\linewidth}\raggedleft
deficiency zero
\end{minipage} & \begin{minipage}[b]{\linewidth}\raggedleft
seven-support only
\end{minipage} & \begin{minipage}[b]{\linewidth}\raggedleft
signed-service only
\end{minipage} & \begin{minipage}[b]{\linewidth}\raggedleft
residual-pair only
\end{minipage} & \begin{minipage}[b]{\linewidth}\raggedleft
pre-tier residual
\end{minipage} \\
\midrule\noalign{}
\endhead
\bottomrule\noalign{}
\endlastfoot
positive & 4,761 & 187 & 110 & 924 & 6 & 2 & 1 & 3,531 \\
signed & 408 & 0 & 0 & 50 & 0 & 0 & 0 & 358 \\
\end{longtable}
}

The exact S-tier-superlevel cut gives:

{\small\def\LTcaptype{none} % do not increment counter
\begin{longtable}[]{@{}
  >{\raggedright\arraybackslash}p{(\linewidth - 10\tabcolsep) * \real{0.1429}}
  >{\raggedleft\arraybackslash}p{(\linewidth - 10\tabcolsep) * \real{0.1905}}
  >{\raggedleft\arraybackslash}p{(\linewidth - 10\tabcolsep) * \real{0.1905}}
  >{\raggedleft\arraybackslash}p{(\linewidth - 10\tabcolsep) * \real{0.1905}}
  >{\raggedright\arraybackslash}p{(\linewidth - 10\tabcolsep) * \real{0.1429}}
  >{\raggedright\arraybackslash}p{(\linewidth - 10\tabcolsep) * \real{0.1429}}@{}}
\toprule\noalign{}
\begin{minipage}[b]{\linewidth}\raggedright
family
\end{minipage} & \begin{minipage}[b]{\linewidth}\raggedleft
pre-tier
\end{minipage} & \begin{minipage}[b]{\linewidth}\raggedleft
tier-certified
\end{minipage} & \begin{minipage}[b]{\linewidth}\raggedleft
residual
\end{minipage} & \begin{minipage}[b]{\linewidth}\raggedright
certified SHA-256
\end{minipage} & \begin{minipage}[b]{\linewidth}\raggedright
residual SHA-256
\end{minipage} \\
\midrule\noalign{}
\endhead
\bottomrule\noalign{}
\endlastfoot
positive & 3,531 & 1,219 & 2,312 &
\texttt{\seqsplit{744d872920309c361d6d7f806f140a696e3fc3ae0f75d760d8a07f304d562b6b}}
&
\texttt{\seqsplit{0297ba35311c757cd5c6ec548d2af18410dfd37e791c7679de932fe4bf38695b}} \\
signed & 358 & 159 & 199 &
\texttt{\seqsplit{7f59ea94fe876205ccb72dc97b026b2954feac62375122634aafa318084428ee}}
&
\texttt{\seqsplit{1a9c06123645855d3b4f23d4886b0ada3c3ff3614fc94a7d22c01f411c1355c8}} \\
\end{longtable}
}

Therefore the exact post-tier residual is

\[
 2312+199=2511.
 \tag{4.1}
\]

Although the legacy sufficient cut falsely rejected 208 individual
pair--descriptor incidences, every affected pair has another genuine
failed descriptor. Consequently the pair sets and fingerprints in this
section are unchanged.

\subsection{5. Corrected affine-feasibility
calculation}

For a primitive descriptor weight \(w\), at each positive level \(r\)
put

\[
 E_r=\{i:w_i=r\},\qquad L_r=\{i:w_i<r\}.
\]

For a rational stoichiometric subspace \(S\), the descriptor is
realizable in some affine class exactly when, for every \(r\), there is
\(v_r\in S\) with

\[
 (v_r)_i=0\ (i\in L_r),\qquad (v_r)_i>0\ (i\in E_r).
 \tag{5.1}
\]

For a specified affine class one also needs a base point having the
exact eventual values \(0\), \(1\), or at least \(2\) in the zero-weight
coordinates.

Necessity of (5.1) follows from Gordan's alternative: failure produces
an affine invariant zero above level \(r\), nonnegative and nonzero at
level \(r\), which contradicts constancy after normalization by that
scale. Sufficiency follows from the polynomial flag

\[
 x(n)=b+\sum_r n^r v_r.
\]

Applying this exact test only to the genuine failures from Theorem 2.1
gives:

{\small\def\LTcaptype{none} % do not increment counter
\begin{longtable}[]{@{}
  >{\raggedright\arraybackslash}p{(\linewidth - 8\tabcolsep) * \real{0.1579}}
  >{\raggedleft\arraybackslash}p{(\linewidth - 8\tabcolsep) * \real{0.2105}}
  >{\raggedleft\arraybackslash}p{(\linewidth - 8\tabcolsep) * \real{0.2105}}
  >{\raggedleft\arraybackslash}p{(\linewidth - 8\tabcolsep) * \real{0.2105}}
  >{\raggedleft\arraybackslash}p{(\linewidth - 8\tabcolsep) * \real{0.2105}}@{}}
\toprule\noalign{}
\begin{minipage}[b]{\linewidth}\raggedright
family
\end{minipage} & \begin{minipage}[b]{\linewidth}\raggedleft
corrected failures
\end{minipage} & \begin{minipage}[b]{\linewidth}\raggedleft
feasible
\end{minipage} & \begin{minipage}[b]{\linewidth}\raggedleft
infeasible
\end{minipage} & \begin{minipage}[b]{\linewidth}\raggedleft
legacy false failures removed
\end{minipage} \\
\midrule\noalign{}
\endhead
\bottomrule\noalign{}
\endlastfoot
positive & 12,250 & 9,349 & 2,901 & 200 \\
signed & 428 & 360 & 68 & 8 \\
\textbf{total} & \textbf{12,678} & \textbf{9,709} & \textbf{2,969} &
\textbf{208} \\
\end{longtable}
}

The 151-pair classwise branch is unchanged:

{\small\def\LTcaptype{none} % do not increment counter
\begin{longtable}[]{@{}
  >{\raggedright\arraybackslash}p{(\linewidth - 8\tabcolsep) * \real{0.1765}}
  >{\raggedleft\arraybackslash}p{(\linewidth - 8\tabcolsep) * \real{0.2353}}
  >{\raggedleft\arraybackslash}p{(\linewidth - 8\tabcolsep) * \real{0.2353}}
  >{\raggedright\arraybackslash}p{(\linewidth - 8\tabcolsep) * \real{0.1765}}
  >{\raggedright\arraybackslash}p{(\linewidth - 8\tabcolsep) * \real{0.1765}}@{}}
\toprule\noalign{}
\begin{minipage}[b]{\linewidth}\raggedright
family
\end{minipage} & \begin{minipage}[b]{\linewidth}\raggedleft
certified
\end{minipage} & \begin{minipage}[b]{\linewidth}\raggedleft
remaining
\end{minipage} & \begin{minipage}[b]{\linewidth}\raggedright
certified SHA-256
\end{minipage} & \begin{minipage}[b]{\linewidth}\raggedright
remaining SHA-256
\end{minipage} \\
\midrule\noalign{}
\endhead
\bottomrule\noalign{}
\endlastfoot
positive & 143 & 2,169 &
\texttt{\seqsplit{f48882aa1ff52c1594a71fd217fa559492c7010e950285a9fa2e60e02b487b76}}
&
\texttt{\seqsplit{6763a44c9c312c440997a054f7966347d101e3236cdef9ecb90599226de10458}} \\
signed & 8 & 191 &
\texttt{\seqsplit{aead73fd44d08789019326cffcd706a776addf0cbc841979a3d54e8c80c5f88d}}
&
\texttt{\seqsplit{f5c7a694bec0241a67b5cf588e1d074c11e00fb9ae3fbc1cee5570f84e9b4483}} \\
\end{longtable}
}

The combined certified-set fingerprint is
\texttt{\seqsplit{55e243945f86d106b920a27e2249a20b7077b5dc718ec06918cca4368e4a6c96}}.

\subsection{6. Class-local entropy-Foster
consequence}

Fix one of the 151 support pairs, arbitrary strongly connected reaction
graphs on its two linkage supports, arbitrary positive rates, and a
closed irreducible population class \(\Gamma\). If \(\Gamma\) is
infinite, it lies in one affine stoichiometric class.

Suppose a tier sequence in \(\Gamma\) failed the top-S descending-source
condition. By Theorem 2.1, its descriptor fails the exact superlevel
cut. Because the sequence lies in one affine class, necessity of (5.1)
makes that failed descriptor affine-feasible. This contradicts the
definition of the 151-pair set. Thus every tier sequence contained in
\(\Gamma\) satisfies

\[
 T^{S,1}_{\{x_n\}}\cap D_{\{x_n\}}\ne\varnothing.
 \tag{6.1}
\]

The cited publication result is David F. Anderson and Jinsu Kim,
\emph{Some Network Conditions for Positive Recurrence of Stochastically
Modeled Reaction Networks}, \textbf{SIAM Journal on Applied Mathematics}
78(5) (2018), 2692--2713,
\href{https://doi.org/10.1137/17M1161427}{doi:10.1137/17M1161427},
\textbf{Theorem 9}. Theorem 4.2 is the numbering only in arXiv v3.

Journal Theorem 9 is stated with an ambient-space premise. Its
contradiction proof restricts as follows. If the entropy drift were not
at most \(-1\) outside a finite subset of \(\Gamma\), properness would
give an unbounded sequence \(x_n\in\Gamma\) with drift greater than
\(-1\). Extracting the tier subsequence and using (6.1) repeats the
proof of Theorem 9 and forces the drift to \(-\infty\), a contradiction.
Closedness makes the restricted generator equal to the original
generator at every state of \(\Gamma\).

The chain is nonexplosive independently: a reaction with a quadratic
source cannot increase total population because its target is also
binary. Every population-increasing reaction therefore has source degree
at most one, total increasing intensity \(O(1+|x|)\), and bounded jump
size. A linear pure-birth comparison bounds total population on finite
time intervals; each population sublevel contains finitely many states
and has bounded total rates. Continuous-time Foster applied to the
closed irreducible class now proves positive recurrence.

\subsection{7. Reproducibility and frozen derivative
bytes}

Run only the symbolic/set-identity derivative:

\begin{CodeBlock}
PYTHONPATH=src python3 src/s_tier_superlevel_interface.py
PYTHONPATH=src python3 -m unittest \
  tests.test_s_tier_superlevel_interface -v
\end{CodeBlock}

The derivative source does not call the legacy orientation-witness
constructor. Its finite certificate hash is

\texttt{\seqsplit{77c7ce0d2325379acfed7b13a44f9577454279275918ee14f968e313b488a7e0}}.

At creation, the exact derivative code bytes were:

{\small\def\LTcaptype{none} % do not increment counter
\begin{longtable}[]{@{}
  >{\raggedright\arraybackslash}p{(\linewidth - 2\tabcolsep) * \real{0.5000}}
  >{\raggedright\arraybackslash}p{(\linewidth - 2\tabcolsep) * \real{0.5000}}@{}}
\toprule\noalign{}
\begin{minipage}[b]{\linewidth}\raggedright
file
\end{minipage} & \begin{minipage}[b]{\linewidth}\raggedright
SHA-256
\end{minipage} \\
\midrule\noalign{}
\endhead
\bottomrule\noalign{}
\endlastfoot
\texttt{src/\texttt{\seqsplit{s\_tier\_superlevel\_interface.py}}} &
\texttt{\seqsplit{1a4e27fcf40af76cac6281f8830b7644bf086b3c05d97a963ce9f5bac736ad57}} \\
\texttt{\seqsplit{tests/\texttt{\seqsplit{test\_s\_tier\_superlevel\_interface.py}}}} &
\texttt{\seqsplit{4d9f960d89a361a27d9dadbd765297783755b9ec09c60f32301229447f51af40}} \\
\end{longtable}
}

The accompanying hash manifest pins this note and the hostile audit
after their final byte pass.
